% Folder 136, batch 06 — archivists' pages 101–120.
% Pass: Opus 5.5 (claude-opus-5-5), 2026-10-03 — first pass, unchecked against the pages by a human.
%
% Skipped as unrelated: pp. 102, 106, 111, 113, 115 (typescript of a text on
% monodromy pairings and Néron models, used as scrap paper, no ink of his on
% the mathematics in hand); p. 108 (English typescript by another hand);
% p. 109 (blank ruled sheet); p. 117 (the opening of a personal letter, no
% mathematics); p. 119 (third-party letter, not circulated). Page 104 is a
% typescript page of the same kind, kept only for the diagram he drew at its
% foot. Page 107 is in very faint pencil; much of it is lost.
\documentclass[11pt,a4paper]{article}
\input{../preamble/grothendieck}

\folder{136}
\batch{6}
\pages{101}{120}
\foldertitle{Complexe de De Rham à puissance divisée [conférence de 1976 à l’IHÉS] : notes manuscrites (s.d.).}
\dating{[à partir de 1975-1976]}
\watermark{Édition de démonstration}

\begin{document}

\page{101}
\note{la page continue une construction commencée avant le lot : les algèbres de De Rham « spéciales » $A^{**}_{*}$, les catégories $\mathrm{DRsp}$, $\mathrm{DRspfl}$ et l'élément $T$ sont fixés dans les pages précédentes. La lettre gothique notée ici $\mathfrak{G}$ (dans $\mathfrak{G}_{*}$, $\mathfrak{G}_{*k}$) est d'une lecture incertaine ; la notation est tenue telle quelle dans tout le lot.}

\struck{Considérons l'\uncertain{hom.} induit en degré $0$}

\struck{$k_0 = k \to A_0$}

On trouve donc un \uncertain{hom.}, donnant lieu à un diag. comm.

\begin{tikzcd}
A_{*} \arrow[r] & \mathfrak{G}_{*} \otimes_{\mathbb{Z}} A_0 \\
k_{*} \arrow[u] \arrow[r, hook] & \mathfrak{G}_{*k} \;(= \mathfrak{G}_{*} \otimes_{\mathbb{Z}} k) \arrow[u, "{\mathfrak{G}_{*} \otimes_{\mathbb{Z}} u_0}"']
\end{tikzcd}

Ainsi on trouve un diagramme canonique d'hom. dans $\mathrm{DRsp}(\mathrm{Ss}, k_{*})$

\begin{tikzcd}
\mathrm{DR}(A^{10}_{*}, T) \arrow[r] \arrow[d] & A^{**}_{*} \\
\mathrm{DR}(A^{10}_{0} \otimes_{\mathbb{Z}} \mathfrak{G}_{*}, T) & \mathrm{DR}(\mathfrak{G}_{*k}, T) \arrow[l]
\end{tikzcd}

reliant une \add{\ill{}} algèbre de DR spéciale, en $3$ crans, à l'algèbre de DR standard $\mathrm{DR}(\mathfrak{G}_{*k}, T)$ (qui semble la plus économique de toutes). Si $A^{1 0}_{*}$ est \add{\ill{}} flasque, alors il en est de même de $\mathrm{DR}(A^{10}_{*}, T)$, et le diagramme ci-dessus se trouve dans $\mathrm{DRspfl}$.
\note{dans le second diagramme, une petite flèche oblique est tracée au-dessus du $T$ de $\mathrm{DR}(\mathfrak{G}_{*k}, T)$.}

On va introduire maintenant les catégories $\mathrm{DRsp}'(\mathcal{X}, k)$, $\mathrm{DRspfl}'(\mathcal{X}, k)$, ayant les mêmes objets que les catégories sans ${}'$, mais ayant plus de morphismes : les hom \struck{\ill{}} \add{différentielles} \struck{\ill{}} à p.d., \uncertain{respectant} $T$

\page{103}
\[
u : \widehat{A}^{**} \to \widehat{A}'^{**},
\qquad
\widehat{A}^{**} = \prod A^{ij},
\quad
\widehat{A}'^{**} = \prod A'^{ij}
\]
hom. d'\uncertain{anneaux} \uncertain{de} $\widehat{A}$, $\widehat{A}'$, \uncertain{appliquant} $\widehat{A}^{**+}$ dans $\widehat{A}'^{**+}$ et compatible aux p.d., et tel que $u(T) = T$ (\uncertain{cela} a priori $u : k\{T\} \to k\{T\}$, où $k\{T\} \simeq H^{*}(A^{**})$, resp. $\simeq H^{*}(A'^{**})$, \uncertain{bien} compatible aux p.d.). \struck{À \ill{}, il}

On désigne \struck{\ill{} \ill{} de \ill{}} par $\mathrm{DRsp}'(\mathcal{X}, k)_0$, $\mathrm{DRspfl}'(\mathcal{X}, k)_0$ la cat. \uncertain{déduite} \uncertain{ainsi} (i.e. les groupoïdes engendrés). On a donc des \struck{\ill{}} diagrammes comm. de foncteurs

\begin{tikzcd}
\mathrm{DRsp}(\mathcal{X}, k) \arrow[r, "\text{surj. sur ob.}"] \arrow[d, "\text{surjectif sur obj.}"'] & \mathrm{DRsp}(\mathcal{X}, k)_0 \\
\mathrm{DRsp}'(\mathcal{X}, k) \arrow[r, "\text{surj. sur objets}"'] & \mathrm{DRsp}'(\mathcal{X}, k)_0 \\
\mathrm{DRspfl}(\mathcal{X}, k) \arrow[r, "\text{surj. sur ob.}"] \arrow[d, "\text{surjectif sur ob.}"'] & \mathrm{DRspfl}(\mathcal{X}, k)_0 \arrow[d] \\
\mathrm{DRspfl}'(\mathcal{X}, k) \arrow[r, "\text{surj. sur ob.}"'] & \mathrm{DRspfl}'(\mathcal{X}, k)_0
\end{tikzcd}

\note{transcription simplifiée du diagramme de la page. Les objets de droite sont suivis de « $\simeq$ \uncertain{pleinement} \ill{} » (première et troisième lignes) et d'un « $\simeq$ » isolé (quatrième ligne) ; une flèche verticale descend aussi de $\mathrm{DRsp}(\mathcal{X}, k)_0$ vers $\mathrm{DRsp}'(\mathcal{X}, k)_0$. Une annotation près de la flèche verticale de droite inférieure est illisible (« (\ill{} pl. \ill{}) »). De fines flèches courbes, tracées plus tard, relient chaque objet de la colonne de gauche à celui de la ligne suivante, et de même à droite ; elles aboutissent à des points marqués d'un pâté.}

qui montre que les deux catégories \add{\ill{}} $\mathrm{DRsp}'(\mathcal{X}, k)_0$ et $\mathrm{DRspfl}'(\mathcal{X}, k)_0$ sont des groupoïdes\add{, en} \uncertain{fait} \uncertain{triviaux} (\uncertain{réunions} d'une famille d'objets avec \uncertain{isom.} \uncertain{transitifs} \uncertain{entre} eux). Donc ils \struck{sont} \uncertain{com}\ill{}

\page{104}
\note{dessin de sa main au pied d'une page dactylographiée sans rapport ; le haut du diagramme passe sous une bande adhésive.}

\[
\begin{array}{ccc}
\mathrm{DR}(A^{1,0}_{*}, T) & \longrightarrow & A^{**} \\
\big\downarrow & & \\
\mathrm{DR}(\mathfrak{G}_{*} \otimes_{\mathbb{Z}} A_0, T) & & \\
\big\uparrow & &
\end{array}
\]

\note{la flèche ascendante vers $\mathrm{DR}(\mathfrak{G}_{*} \otimes_{\mathbb{Z}} A_0, T)$ n'a pas de source écrite ; le $\mathfrak{G}_{*}$ est récrit par-dessus un autre symbole.}

\page{105}
\note{la première ligne commence au ras du bord ; elle poursuit la phrase laissée à la fin de la p. 103.}

\ill{} équivalent au groupoïde connexe défini par deux groupes, que nous nous proposons de déterminer [dans le cas $\mathcal{X} = (\mathrm{Ss})$].

\begin{tikzcd}
\widehat{A}^{**}_{*} \arrow[r, "u"] & \widehat{A}'^{**}_{*} \\
\widehat{\mathrm{DR}}(A^{10}_{*}, T) \arrow[u] \arrow[ur, dashed]
\end{tikzcd}

\note{la flèche oblique, en pointillé, porte une marque qui ressemble à un « $5$ », peut-être un « $\exists$ ». Sous $\widehat{\mathrm{DR}}(A^{10}_{*}, T)$ : « $= (A_{*})$ », et un double trait oblique renvoie à la formule suivante.}

\[
\Gamma^{*}(A_{**}) \mathbin{\widehat{\otimes}_k} \widehat{\Lambda}^{*}(B_{**}),
\qquad
B_{**} = A_{*}/k_{*}
\]

\[
A_{*} \xrightarrow{\;u\;} A'^{**+},
\qquad
T \longmapsto T
\]

\[
v(p,q) : A_{*} \to A'^{pq}
\qquad (p, q \geqslant 0,\ p + q > 0)
\]
\[
v(1,0)(T) = T,
\qquad
v(p,q)(T) = 0 \ \text{si}\ (p,q) \neq (1,0)
\]
\[
\bar{v}(p,q) : \mathfrak{B} \to A'^{pq}
\qquad (\text{si}\ (p,q) \neq (1,0))
\]

\begin{tikzcd}
k \arrow[r, hook] \arrow[d, no head, "\Vert" description] & A \arrow[d, "{v(1,0)}"] \\
k \arrow[r, hook] & A'^{1,0}
\end{tikzcd}

$k \subset$ \struck{$A$}

\begin{tikzcd}
 & & \mathfrak{G}_{*} \otimes A \\
k \arrow[r, hook] & A \arrow[r] \arrow[ur, no head] \arrow[dr] & \mathfrak{B} \arrow[d] \\
 & & A'^{pq}
\end{tikzcd}

\note{sous la ligne $k \subset A$, un trait relie aussi $k$ à $A'^{pq}$ ; à droite de $A'^{pq}$ un signe « $=$ » sans suite.}

\page{107}
\note{page au crayon, très pâle ; seuls des fragments se lisent, et l'ordre des lignes est celui de la page. Le texte est compressé.}

Détermination des hom. \struck{dans}
$\mathrm{DRsp}$ de $\mathrm{DR}^{**}_{*k} = \Gamma_k(\mathfrak{G}_{*k}) \mathbin{\widehat{\otimes}} \widehat{\Lambda}_k(\mathfrak{G}_{*k}/k)$ dans une \uncertain{algèbre} de De Rham \ill{} \ill{} \ill{} \uncertain{augmentées}, \ill{} \ill{} \ill{} différentielle, et \ill{} $T$.
\ill{} \ill{} $\mathfrak{G}_{*k} \to \widehat{A}^{**+}_{*}$ \ill{}
\ill{} \ill{} $(T) = T$)
\ill{}
On \ill{} $\varphi$ \ill{} $\widehat{A}^{**+}$ \ill{} \ill{} \ill{}
\ill{} $\varphi : \mathfrak{G}_{*k} : k \ast k \to \widehat{A}^{**+}_{*}$ \ill{} \ill{}

\[
\xi_0 = \beta_1(\ldots),
\qquad
\xi_0 \in \widehat{A}^{**+}_{1}
\]
\uncertain{il} \uncertain{compatible} avec les deux augmentations
$\alpha^{*}, \beta^{*} : \widehat{A}^{**+}_{1} \to \widehat{A}^{**+}_{0}$
\[
\boxed{\alpha^{*}(\xi_0) = 1, \quad \beta^{*}(\xi_0) = 1}
\qquad (\alpha, \beta : \Delta_0 \to \Delta_1)
\]
les conditions \ill{} \ill{} s'expliquent \ill{}
\[
\boxed{\xi_0 \, d\xi_0 = d\xi_0 \cdot \xi_0}
\qquad
\boxed{d\xi_0 \cdot \xi_0 \ \ldots}
\]
\ill{} \ill{} $\Delta_2 \to \Delta_1$, et \uncertain{si} $q_0$, $q_1$ \ill{} les \ill{}
$(q_0(0) = 0,\ q_0(1) = q_1(0) = 1,\ q_1(\ldots) \ldots)$, puis
\[
q_0^{*}(\xi_0)\, q_1^{*}(\xi_0) = q_1^{*}(\xi_0)\, q_0^{*}(\xi_0)
\]
\[
q_0^{*}(\xi_0)\, q_1^{*}(d\xi_0) \ldots = \ldots\, q_0^{*}(d\xi_0)\, q_1^{*}(\xi_0)
\]
\[
q_1^{*}(\xi_0)\, q_0^{*}(d\xi_0) \ldots = \ldots\, q_0^{*}(\xi_0)\, q_1^{*}(d\xi_0)
\]
\[
q_0^{*}(d\xi_0)\, q_1^{*}(d\xi_0) = - q_1^{*}(d\xi_0)\, q_0^{*}(d\xi_0)
\]
\note{les quatre relations sont encadrées ensemble ; dans les deux du milieu des coefficients écrits en surcharge sont illisibles, d'où les « $\ldots$ ».}

\section{Foncteurs de De Rham simpliciaux}
\note{titre de sa main, ajouté en tête de la p. 110 ; « simpliciaux » est une lecture douteuse.}

\page{110}
On s'intéresse aux foncteurs
\[
C^{**} : (\mathrm{Ss})^{\circ} \longrightarrow \mathfrak{D}
\]
où $\mathfrak{D}$ est la catégorie des $S_k$-algèbres bigraduées \add{alternées} différentielles \uncertain{augmentées}, à p.d. (bidegrés positifs, diff. de degré $+1$, …).

et on s'intéresse à ceux qui transforment \emph{lim} quelconques de $(\mathrm{Ss})$ en \emph{lim}, donc qui sont représentables. Ils \struck{donnent} \add{peuvent} s'interpréter \add{(\uncertain{fantasmes})} \struck{comme} des $S_k$-algèbres bigraduées alternées … ${}^{1\!}/_{2}$ simpliciales, i.e. des foncteurs
\[
\Delta^{\circ} \longrightarrow \mathfrak{D}
\]
et on \uncertain{interprète} dans les \uncertain{notations} les deux sortes d'objets, écrivant
\[
C^{**} = C^{**}_{*} = (C^{**}_{n})_{n \geqslant 0} \ \ldots
\]
et de même pour les morphismes entre foncteurs, correspondant aux morphismes entre ensembles ${}^{1\!}/_{2}$ simpliciaux, respectant toutes les structures.

À $C^{**}_{*}$ est associé un foncteur « cohomologique de De Rham associé à $C^{**}_{*}$ »
\[
\mathcal{H}^{**}_{C}(\mathcal{X}) = H^{**}(C^{**}(\mathcal{X}))
\]
à valeurs dans les $S_k$-algèbres bigraduées alternées. On s'intéresse aux $C^{**}$ \struck{$(\mathcal{X})$} (dits admissibles) pour lesquels pour tout $\mathcal{X}$,
\note{le $H^{**}$ du second membre porte un indice récrit, peut-être $\varepsilon$.}

\page{112}
\[
\tau_{\omega}(H^{**}(\mathcal{X}))
\]
[qui est une algèbre graduée \struck{\ill{}} alternée, entre autres, de $M_{\omega}$] \struck{provient d'un} \uncertain{est} dans l'image essentielle de $\mathrm{Mod}(k)$, i.e. peut s'écrire (de façon \ill{} \ill{} \ill{} \uncertain{unique} \ill{})
\[
\mathcal{H}^{**}_{C}(\mathcal{X}) = H^{*}_{c}(X) \otimes_k S_k
\]
[Rappel. — on sait que $H^{*}_{c}(X)$ \uncertain{hérite} dans la structure de l'algèbre graduée de $\tau_{\omega}(H^{**}(\mathcal{X}))$. Cela implique que pour les $\mathcal{C}^{**}_{n}$ \ill{} \uncertain{la} \uncertain{propriété} (faire $\mathcal{X} = \Delta_n$).
\note{dans la formule centrale, le $H^{*}_{c}$ du second membre est écrit par-dessus un autre $H$.}

\struck{plus que \add{$\forall \mathcal{X}$,} l'augmentation $S_k \to C^{**}(\mathcal{X})$ définie}
\struck{On veut de} On va préciser cette exigence.

On peut considérer $\mathcal{C}^{**}_{n}$ comme un faisceau (à structures de type $\mathfrak{D}$) sur le topos $(\mathrm{SS})$, et on a l'augmentation par le faisceau constant $S_k$
\[
S_k \to C^{**}_{n} .
\]
\ill{} NB, on regarde ici la \uncertain{catégorie} \uncertain{des} \ill{} comme un \ill{} \uncertain{principal}] de \uncertain{sorte} que \struck{\ill{}} \add{$\mathcal{D}(S_k)$} \ill{} (des \struck{\ill{}} cat. \uncertain{dérivées} \struck{\ill{}}) \add{« $\mathcal{D}$($S_k$-mod gradués) »} :
\[
\begin{array}{ccc}
\mathbb{R}^{*}\Gamma(\mathcal{X}, S_k) & \longrightarrow & \mathbb{R}^{*}\Gamma(\mathcal{X}, C^{**}) \\
 & & \big\uparrow \\
(1) & & C^{**}(\mathcal{X}) = \Gamma(\mathcal{X}, C^{**})
\end{array}
\]

\page{114}
En fait, \struck{en demandant que} \ill{} \uncertain{appliquant} $\tau_{\omega}$ \uncertain{à} $\mathbb{R}^{*}\tau_{\omega}$, on trouve
\[
\begin{array}{ccc}
\tau_{\omega}(\mathbb{R}^{*}\Gamma(\mathcal{X}, S_k)) & \longrightarrow & \tau_{\omega}\mathbb{R}^{*}\Gamma(\mathcal{X}, C^{**}) \\
 & & \big\uparrow \\
(2) & & \tau_{\omega} C^{**}(\mathcal{X})
\end{array}
\]
\note{deux symboles sont biffés avant les deux membres de la flèche horizontale.}

On \struck{demande} \add{voudrait} que ces deux flèches \struck{soient des isomorphismes} \add{dans la cat. dérivée}, ce qui signifie aussi que les homomorphismes déduits de (1) en passant à la coh.
\[
\begin{array}{ccc}
H^{**}(\mathcal{X}, S_k) & \longrightarrow & H^{**}(\mathcal{X}, C^{**}) \\
 & & \big\uparrow \\
(1') & & H^{**}_{C}(X)
\end{array}
\]
soient à noyau et conoyau négligeables, ce qui \struck{signifie} \add{implique} \uncertain{précisément}, pour la flèche horizontale, que
\[
(2') \quad
\begin{array}{ccc}
\tau_{\omega}(H^{*}(\mathcal{X}, k) \otimes_k S_k) & \xrightarrow{\ \sim\ } & \tau_{\omega} H^{**}(\mathcal{X}, C^{**}) \\
 & & \big\uparrow \wr \\
 & & \tau_{\omega}(H^{**}_{c}(X))
\end{array}
\]
\note{une accolade embrasse la formule (2') ; une ligne pointillée descend en diagonale de la flèche horizontale vers $\tau_{\omega}(H^{**}_{c}(X))$. Le premier $H^{*}$ est récrit par-dessus un autre symbole.}

\page{116}
Prenons $\mathcal{X} = \Delta_n$, alors \struck{le} \add{la} diagramme (1) devient
\[
\begin{array}{ccc}
\mathbb{R}^{*}\Gamma(\Delta_n, S_k) \simeq S_k & \longrightarrow & \mathbb{R}^{*}\Gamma(\Delta_n, C^{**}) = \mathcal{C}^{**}_{n} \\
 & & \big\uparrow \ \mathrm{id}_{C^{**}_{n}} \\
 & & C^{**}(\Delta_n) = C^{**}_{n}
\end{array}
\]
et la condition énoncée se réduit à celle concernant la flèche horizontale, qui signifie aussi
\[
(3) \qquad
H^{**}_{C}(\Delta_n) = H^{**}(C^{**}_{n})
\quad
\text{\struck{est isom.}}
\]
\[
\uparrow \varepsilon \ \text{augmentation}
\qquad
S_k
\]
l'augmentation est un quasi-isom., i.e.
\[
\text{3a)} \qquad
\text{\struck{$\ldots$}}\ S^{i}_{k} \to H^{i0}(C^{**}_{n})
= \mathrm{Ker}(C^{i,0}_{n} \to C^{i1}_{n})
\]
\ill{} \ill{}, iso pour $i$ grand ($i \geqslant \lambda(n)$ disons)
\[
\text{3b)} \qquad
\forall p > 0, \quad H^{ip}(C^{**}_{n}) = 0
\ \text{pour}\ i\ \text{grand (dépendant de}\ p\text{)}
\]
soit
\[
H^{ip}(C^{**}_{n}) = 0 \quad \text{si} \quad i \geqslant \lambda(p)
\]
\note{$\lambda(p)$ est récrit sur une autre lettre biffée.}
où on peut supposer $\lambda : \mathbb{N}^{*} \to \mathbb{N}^{*}$ strictement croissante.

A priori, on $\lambda = \lambda_n$ dépend du choix de $n$. On suspecte qu'il faut pouvoir prendre $\lambda$ indépendant de $n$ …

\page{118}
dans ce cas la flèche horizontale sera bien un isomorphisme \struck{donc} après application de $\tau_{\omega}$… En effet, chaque $C^{i*}_{*}$ est une résolution de $(S^{i}_{k})_{*}$ (dans $(\mathrm{SS})$), jusqu'au degré $p$ si $i$ est assez grand par rapport à $p$, donc ….

\struck{On} Regardons la condition d'isomorphie (\uncertain{application} de 2) pour la flèche verticale, elle \struck{est \ill{}} est conséquence de la condition suivante (et sans doute équivalente)

(4) Pour tout $i \geqslant 0$ considérons les $C^{ip}_{*}$ \add{\uncertain{gradués} \ill{}} $(p \geqslant 0)$ et soit $\pi(i)$ le plus petit entier (s'il existe) tel que $C^{ip}_{*}$ soit flasque pour $p \geqslant \pi(i)$ ($\pi(i) = +\infty$ sinon), \struck{\ill{}} \add{(\uncertain{pour} $\partial$)} i.e. que ses \uncertain{faisceaux} \ill{} simpliciaux \add{\uncertain{soient}} \ill{}. On veut que $\pi(i) \to +\infty$ quand $i \to +\infty$.

\marginal{(en oblique dans la marge gauche) On veut \struck{\ill{}} affaiblir (4) en \uncertain{(4)'}, en \uncertain{remplaçant} « flasque » par « \ill{} \uncertain{générateurs} \ill{} » (\ill{} $\mathrm{SS}$)}

On aimerait prouver un th. d'unicité (\ill{} déjà l'existence …) Les $C^{**}_{*}$ ayant les structures précisées (au \uncertain{besoin} \ill{} et dont la structure p.d.) et satisfaisant

\page{120}
aux conditions supplémentaires \uncertain{techniques} \uncertain{que} \uncertain{énoncées} (être une « quasi » résolution de $S_k$, et être « quasi-flasque » …)

On note que les hom dans \struck{\ill{}} \add{la} catégorie des objets \add{$C^{**}$} admissibles, sont déjà des « quasi-isom. » (p.r. à $\tau_{\omega}$) — il y a lieu de passer à la cat. des fractions en \struck{les} rendant inversibles tous ces hom., i.e. le groupoïde engendré par la catégorie envisagée. On aimerait donc prouver que ce groupoïde est connexe et simplement connexe ? [NB rappelons \ill{} \add{\ill{}} $S_k = k\{T\}$ \ill{} \ill{} \ill{} plus \uncertain{précis} \uncertain{qu'il} y ait \ill{} \ill{} \uncertain{graduées}.]
\note{le crochet ouvert par « [NB » se ferme après une ligne et demie d'insertions serrées, dont l'ordre est incertain.}

Cela signifie donc

\uncertain{Connexité} : si $C^{**}_{*}$ et $C'^{**}_{*}$ sont des théories admissibles, $\exists$ théorie admissible $C''^{**}_{*}$ et hom
\[
C''^{**}_{*} \to C^{**}_{*},
\qquad
C''^{**}_{*} \to C'^{**}_{*} .
\]

b) Si on a deux hom
\[
C'^{**}_{*} \rightrightarrows C^{**}_{*}
\]
$\exists$ hom
\[
C''^{**}_{*} \to C'^{**}_{*}
\]
qui « égalise ».

\marginal{(un grand « ? » à gauche, en regard des deux conditions ; à droite, séparé par un trait oblique) il y aurait plus de chances avec catégorie de fractions à droite ???}

\note{la page s'arrête sur ces deux conditions ; l'argument se poursuit au-delà du lot.}

\end{document}
